\section{Modal Geometry of NDH}
\label{sec:ndh-modal-geometry}

NDH is structured around three modal regions $M_A, M_B, M_C$ forming a triadic
partition of the underlying manifold $M$. Each modal region carries its own
geometric signature, curvature profile, and adjacency behavior. Modal geometry
defines how NDH transitions between cognitive modes through geometric traversal.

\subsection{Modal Regions}

Each modal region $M_i$ is an open, connected submanifold of $M$:
\[
M_i \subset M, \quad i \in \{A, B, C\}.
\]
The modal regions intersect along modal boundaries:
\[
\partial M_i \cap \partial M_j \neq \emptyset
\]
for all $i \neq j$. These intersections define modal adjacency.

\subsection{Modal Boundaries}

Modal boundaries are codimension-one submanifolds:
\[
\mathcal{B}_{ij} = \partial M_i \cap \partial M_j.
\]
Traversal across $\mathcal{B}_{ij}$ induces a modal transition. The geometry of
$\mathcal{B}_{ij}$ determines the continuity of spectral emission and recursion
behavior across modes.

\subsection{Modal Curvature}

Each modal region carries a curvature function:
\[
\kappa_i : M_i \to \mathbb{R}.
\]
Modal curvature governs the geometric behavior of recursion operators and
spectral components. Modal transitions preserve curvature up to a bounded
distortion:
\[
|\kappa_i(x) - \kappa_j(y)| \le \epsilon
\]
for $x \in M_i$, $y \in M_j$ with $x$ and $y$ adjacent across $\mathcal{B}_{ij}$.

\subsection{Modal Adjacency Graph}

Modal adjacency is encoded by a graph:
\[
\mathcal{G} = (V, E)
\]
with
